% ============================================================================
% Appendix A — Implementation reference
% Body prose in the chapters is concept-only (no module paths,
% no class or function names, no script names, no folder or file references).
% All of those references that a re-implementer needs live here.
% ============================================================================

\chapter{Implementation reference}
\label{app:implementation}

\section{Code and data availability}
\label{app:code-data-availability}

The implementation of the architecture, the cycle runner, the calibration scripts, the baseline drivers, the verifier, the retrieval-utility classifier and the table and figure generators is released at \url{https://github.com/aksaN000/caem-thesis}. The results are in the shared Google Drive folder \url{https://drive.google.com/drive/folders/1xErhsRDYe_qO0ZgUBn0WXK81eHc5jZCa}: \texttt{1\_thesis/} holds this report, \texttt{2\_results\_phase1e/} every number it reports (\texttt{caem/} for the main run, \texttt{baselines/} for the final baseline panel, \texttt{analysis/} for the scored summaries and tests), \texttt{3\_ruc\_ablation\_phase1d/} the trajectory without the classifier, and \texttt{9\_archive/} superseded files, kept rather than deleted. A \texttt{README} in each folder maps it to the repository's output directories. The $64$-gigabyte Wikipedia passage index and the pre-main snapshot are mirrored on the Hugging Face dataset repository \texttt{aksaN000/caem-passage-index-21m} (private; access on request; \Cref{app:repro-snapshot}).

\section{Repository layout and module map}
\label{app:repo-layout}

Body prose throughout the thesis refers to the architectural components by concept name; the inventory below maps each concept to the file that implements it. Paths are relative to the repository root.

{\scriptsize
\setlength{\tabcolsep}{4pt}
\renewcommand{\arraystretch}{1.15}
\begin{longtable}{@{}>{\ttfamily}p{5.8cm}p{2.2cm}p{6.6cm}@{}}
\caption[Repository layout]{Repository layout. Body prose throughout the thesis refers to architectural components by their concept names (memory store, three-tier router, verifier ensemble, retrieval-utility classifier, self-improvement loop, retention guard); this table maps each concept to the file that implements it.}
\label{tab:app-repo-layout}\\
\toprule
\textnormal{\textbf{Path}} & \textbf{Role} & \textbf{Description} \\
\midrule
\endfirsthead
\toprule
\textnormal{\textbf{Path}} & \textbf{Role} & \textbf{Description} \\
\midrule
\endhead
\bottomrule
\endfoot
\multicolumn{3}{@{}l}{\textnormal{\textit{Top-level directories}}} \\
caem/              & package root & Core Python package: pipeline, memory, routing, verification, training. \\
scripts/           & runners      & Cycle runner, calibration scripts, baseline runners, rescore utilities. \\
eval/              & harness      & Shared per-benchmark eval scaffold (EM, em\_llm\_judged, CHM, CES). \\
outputs/           & artefacts    & Run outputs: full\_run/ (main run), cycle\_0/, rerun\_v2/ and baselines\_phase1e/ (baselines), analysis/ (scored summaries). \\
data/              & external     & Passage index, calibration pairs, alias dictionary, retention and RUC datasets. \\
tests/             & test suite   & Pytest suite covering storage gate, deferred buffer, retention probe, SIL, verifier. \\
docs/              & runbook      & PRODUCTION\_RUNBOOK.md: deployment-mode spec and cycle-trigger policies. \\
defense/thesis\_report/ & thesis   & LaTeX chapters, bibliography, figures, main.pdf build target. \\
\midrule
\multicolumn{3}{@{}l}{\textnormal{\textit{Pipeline entry points}}} \\
caem/pipeline.py            & per-query  & Pipeline class: serial driver used in production and demo paths. \\
caem/pipeline\_batch.py     & per-cycle  & BatchPipeline class: batched cycle driver used by the experiment runner. \\
scripts/run\_experiment.py  & runner     & Canonical cycle driver: invokes BatchPipeline, SIL, retention probe, retroverify. \\
\midrule
\multicolumn{3}{@{}l}{\textnormal{\textit{Configuration}}} \\
caem/config.py                       & registry   & Central CAEMConfig dataclass; every threshold, weight, structural toggle. \\
caem/ruc/v2/feature\_spec.json       & RUC config & RUC threshold, regulariser, feature order, model paths. \\
\midrule
\multicolumn{3}{@{}l}{\textnormal{\textit{Memory module}}} \\
caem/memory/store.py     & storage gate    & EpisodicMemoryStore: fixed-threshold admission and FAISS-backed retrieval. \\
caem/memory/deferred.py  & deferred buffer & DeferredBuffer and DeferredEntry: TTL-bounded reconsider queue. \\
caem/memory/encoder.py   & embeddings      & SBERT-side encoder for memory keys and Tier 1 lookup. \\
\midrule
\multicolumn{3}{@{}l}{\textnormal{\textit{Routing}}} \\
caem/routing/router.py & tier dispatch & Tier 1/2/3 router with safety override and combined-score thresholds. \\
caem/routing/ruc.py    & RUC head      & LRRetrievalUtilityClassifier: binary head for Tier 2 versus Tier 3 routing. \\
\midrule
\multicolumn{3}{@{}l}{\textnormal{\textit{Verification}}} \\
caem/verification/verifier.py            & verifier ensemble & UnifiedVerifier: nine-signal composite producing $u_{\text{pre}}$ and $u_{\text{stored}}$. \\
caem/verification/adaptive\_nli\_judge.py & NLI dispatch     & AdaptiveNLIJudge: length-gated dispatcher; in the deployed configuration every hypothesis routes to MiniCheck under a 512-token truncation contract because the long-hypothesis Qwen-3B backend was ablated (\Cref{sec:adj-judge-ablation}). \\
caem/verification/cal\_prob\_composite.py & calibration      & Per-signal isotonic curves plus per-benchmark shrinkage and boost layer. \\
caem/retrieval/rag.py                    & retrieval         & FAISS-backed passage index reader for Tier 3 grounded generation. \\
\midrule
\multicolumn{3}{@{}l}{\textnormal{\textit{Training}}} \\
caem/training/self\_improvement.py & SIL trainer     & SelfImprovementLoop: LoRA fine-tune and retroverify hook. \\
caem/training/retention\_probe.py  & retention guard & Per-cycle multi-modal probe panel and ratio-floor check. \\
caem/training/pool\_reweighting.py & pool balancer   & Per-benchmark temperature-mixed reweighting and share ceiling. \\
caem/training/loop\_filter.py      & SIL admission   & Per-pair filter at $\tau_{\text{train}} = 0.70$ above the storage cut-point. \\
\midrule
\multicolumn{3}{@{}l}{\textnormal{\textit{Evaluation}}} \\
eval/harness.py            & eval scaffold     & Per-benchmark EM, F1, latency, CHM, and CES. \\
eval/benchmarks.py         & benchmark loaders & TriviaQA, FEVER, CommonsenseQA, TruthfulQA, StrategyQA, MMLU loaders. \\
caem/eval/equilibrium.py   & equilibrium       & Equilibrium summary and the automatic early-stop check. \\
\midrule
\multicolumn{3}{@{}l}{\textnormal{\textit{Final results}}} \\
scripts/final\_sources.py      & sources   & The one answer file per system and benchmark that every analysis reads. \\
scripts/make\_final\_results.sh & results   & Rebuilds every number, table and figure of \Cref{ch:evaluation} from the answer files. \\
\end{longtable}
}

\section{Configuration registry}
\label{app:config-registry}

Every constant that governs runtime behaviour lives in the \verb|CAEMConfig| dataclass at \verb|caem/config.py|. The exception is the retrieval-utility classifier, whose configuration lives in \verb|caem/ruc/v2/feature_spec.json| at deployment time (the classifier is loaded once at pipeline construction and is not refit at cycle boundaries; see \Cref{sec:ruc}). The registry below groups every field by concept, records the deployed value, and tags the provenance as \emph{literature} (anchored to a cited paper), \emph{design} (architectural decision documented before any run), \emph{calibrated} (selected empirically at cycle zero), or \emph{pre-specified} (chosen before any data is observed).

{\footnotesize
\setlength{\tabcolsep}{4pt}
\renewcommand{\arraystretch}{1.1}
\begin{longtable}{@{}p{4.6cm}p{2.0cm}p{1.8cm}p{5.8cm}@{}}
\caption[Configuration registry]{Configuration registry. Provenance tags: \emph{literature} (cited paper), \emph{design} (architectural choice documented before any run), \emph{calibrated} (chosen empirically at cycle zero or refit at a cycle boundary), \emph{pre-specified} (registered horizon before any data is observed).}
\label{tab:app-config-registry}\\
\toprule
\textbf{Field} & \textbf{Value} & \textbf{Provenance} & \textbf{Source / rationale} \\
\midrule
\endfirsthead
\toprule
\textbf{Field} & \textbf{Value} & \textbf{Provenance} & \textbf{Source / rationale} \\
\midrule
\endhead
\bottomrule
\endfoot
\multicolumn{4}{@{}l}{\textit{Routing thresholds (\Cref{sec:router})}} \\
\texttt{routing\_lambda}             & $0.70$ & design       & Combined-score weight on similarity. \\
\texttt{tier1\_combined\_threshold}   & $0.90$ & design       & Tier 1 dispatch floor on combined score. \\
\texttt{tier2\_similarity\_threshold} & $0.75$ & design       & Tier 2 dispatch floor on similarity alone. \\
\texttt{safety\_u\_pre\_min}           & $0.38$ & calibrated   & Safety-override floor; lowered from $0.60$ before the main run after an earlier run's first-cycle routing. \\
\texttt{u\_pre\_token\_weight}         & $0.60$ & design       & Token-probability weight in $u_{\text{pre}}$. \\
\texttt{u\_pre\_cconv\_weight}         & $0.40$ & design       & Layer-convergence weight in $u_{\text{pre}}$. \\
\midrule
\multicolumn{4}{@{}l}{\textit{Storage gate thresholds (\Cref{sec:decision})}} \\
\texttt{store\_threshold}            & $0.60$ & design       & STORE if $u_{\text{stored}} \ge$ this. \\
\texttt{defer\_threshold}            & $0.45$ & design       & DEFER if \texttt{defer\_threshold} $\le u_{\text{stored}} < $ \texttt{store\_threshold}. \\
\texttt{abstain\_pground\_ceiling}    & $0.20$ & design       & ABSTAIN if $u_{\text{stored}} < $ \texttt{defer\_threshold} and $p_{\text{ground,max}} <$ this. \\
\texttt{early\_exit\_u\_internal}      & $0.70$ & design       & Confabulation early-exit on internal-uncertainty floor. \\
\texttt{early\_exit\_p\_ground\_max}    & $0.20$ & design       & Confabulation early-exit on grounding ceiling. \\
\midrule
\multicolumn{4}{@{}l}{\textit{Calibration parameters (\Cref{sec:composite})}} \\
\texttt{composite\_mode}             & \texttt{"auto"}  & design     & Auto-switch to calibrated probability composite when the calibration JSON loads. \\
\texttt{composite\_shrinkage\_alpha}  & $0.6$  & calibrated   & Per-benchmark child curve shrinkage toward pooled fit. \\
\texttt{adaptive\_thresholds\_\allowbreak per\_cycle} & \texttt{true} & design   & Refit temperature, isotonic curves, and boost weights at every successful cycle. \\
\texttt{adaptive\_thresholds\_\allowbreak ema\_alpha} & $0.7$ & design       & Exponential-moving-average smoothing on the per-cycle temperature refit. \\
\midrule
\multicolumn{4}{@{}l}{\textit{Retention guard (\Cref{sec:retention})}} \\
\texttt{forgetting\_tolerance}       & $0.93$  & design       & Worst-probe ratio floor; trigger asymmetric rollback if any probe drops below. \\
\texttt{retention\_probes}           & 3-probe panel & design & Multi-modal probe panel \texttt{(mmlu,\allowbreak{} triviaqa\_test,\allowbreak{} commonsense\_qa\_test)}; replaces the single-MMLU probe of v1. \\
\texttt{retention\_probe\_n}          & $200$  & design       & Samples per probe at every cycle. \\
\midrule
\multicolumn{4}{@{}l}{\textit{Deferred buffer (\Cref{sec:deferred-reconsider})}} \\
\texttt{deferred\_buffer\_max\_size}   & $10\,000$ & design     & Bounded FIFO capacity for deferred entries. \\
\texttt{deferred\_buffer\_ttl\_cycles} & $2$     & design       & Reconsider-invocation TTL; entries dropped after two passes. \\
\midrule
\multicolumn{4}{@{}l}{\textit{Retroverify and consolidation (\Cref{sec:retroverify})}} \\
\texttt{retroverify\_prune\_\allowbreak threshold}        & $0.50$ & design & Prune from memory if recalibrated $u_{\text{stored}}$ drops below. \\
\texttt{hit\_counter\_force\_\allowbreak retroverify}      & $10$   & design & Per-entry retrieval-hit threshold forcing reverification next cycle. \\
\texttt{enable\_consolidation}               & \texttt{true} & design & Cycle-boundary SBERT-similarity clustering pass. \\
\texttt{consolidation\_similarity\_\allowbreak threshold} & $0.92$ & design & Minimum cosine similarity for cluster admission. \\
\texttt{consolidation\_search\_k}             & $20$   & design & Per-entry FAISS search depth for neighbours. \\
\texttt{consolidation\_max\_u\_spread}         & $0.15$ & design & Spread guardrail; reject clusters whose $u_{\text{stored}}$ range exceeds this. \\
\midrule
\multicolumn{4}{@{}l}{\textit{Self-improvement training (\Cref{sec:sil})}} \\
\texttt{use\_lora\_training}          & \texttt{true}  & design       & LoRA SIL primary path; replaces the v1 full-fine-tune fallback. \\
\texttt{lora\_r}                     & $32$         & design       & Adapter rank. \\
\texttt{lora\_alpha}                 & $64$         & design       & Scaling factor, set to twice the rank. \\
\texttt{lora\_dropout}               & $0.05$       & design       & Adapter dropout. \\
\texttt{lora\_target\_modules}        & \texttt{(q,k,v,o,\allowbreak{}gate,up,down)\_proj} & literature & All linear projections of the transformer block. \\
\texttt{lora\_learning\_rate}         & $2\!\times\!10^{-4}$ & design & Adapter learning rate. \\
\texttt{use\_8bit\_adamw}             & \texttt{true}  & design       & Applies only to the full fine-tuning fallback; the LoRA path overrides it and uses 32-bit AdamW. \\
\texttt{batch\_size}                 & $4$          & design       & Per-step batch, set for graphics memory. \\
\texttt{grad\_accum\_steps}           & $4$          & design       & Holds the effective batch at $16$. \\
\texttt{epochs\_per\_cycle}           & $3$          & design       & SIL training duration per successful cycle. \\
\texttt{warmup\_steps}               & $500$        & design       & Linear warmup before the standard schedule. \\
\texttt{min\_u\_stored\_for\_training}  & $0.70$       & calibrated   & $\tau_{\text{train}}$; lowered from $0.75$ before the main run. \\
\texttt{pool\_reweighting\_enabled}   & \texttt{true}  & design       & Temperature-mixed reweighting across training benchmarks. \\
\texttt{pool\_reweighting\_\allowbreak temperature} & $2.0$      & design       & Mixture temperature across training benchmarks. \\
\texttt{pool\_reweighting\_\allowbreak upsample\_cap} & $3.0$     & calibrated   & Upper bound on per-benchmark upsampling factor. \\
\texttt{pool\_reweighting\_max\_share} & $0.40$       & design       & Per-benchmark hard share ceiling. \\
\texttt{loop\_distinct4\_threshold}   & $0.25$       & calibrated   & SIL-pool loop filter on token-4-gram distinct ratio. \\
\midrule
\multicolumn{4}{@{}l}{\textit{Verifier ensemble (\Cref{sec:verifier})}} \\
\texttt{verifier\_backend}           & \texttt{"minicheck"} & literature & Short-hypothesis backend; Tang et al.\ 2024. \\
\texttt{minicheck\_model}            & \texttt{lytang/\allowbreak MiniCheck-\allowbreak Flan-T5-\allowbreak Large} & literature & Verifier checkpoint. \\
\texttt{cross\_encoder\_model}        & \texttt{BAAI/\allowbreak bge-reranker-\allowbreak v2-m3} & design & Shared reranker; consumed by grounding and question-answer match. \\
\texttt{verifier\_retrieve\_k}        & $20$  & design & Top-$k$ passages retrieved before reranking. \\
\texttt{verifier\_rerank\_k}          & $3$   & design & Top-$N$ retained after reranking. \\
\texttt{atomic\_min\_tokens}          & $12$  & calibrated & Length gate below which atomic decomposition is bypassed. \\
\texttt{mc\_dropout\_k},\allowbreak{} \texttt{mc\_dropout\_rate} & $5$, $0.1$ & literature & Gal and Ghahramani 2016. \\
\texttt{sc\_chains\_m}                & $3$   & literature & Wang et al.\ 2022 self-consistency. \\
\texttt{se\_samples\_k},\allowbreak{} \texttt{se\_temperature} & $10$, $1.0$ & literature & Farquhar et al.\ 2024 semantic-entropy protocol. \\
\texttt{disable\_h\_norm}             & \texttt{true} & calibrated & Semantic-entropy signal dropped at cycle zero: learned weight near zero, largest verifier cost. \\
\texttt{disable\_p\_contra}           & \texttt{true} & design & Structurally zero under the MiniCheck binary judge. \\
\midrule
\multicolumn{4}{@{}l}{\textit{Generation and retrieval}} \\
\texttt{base\_model\_name}            & \texttt{Qwen/\allowbreak Qwen2.5-3B-\allowbreak Instruct} & design & Apache-licensed three-billion-parameter instruction-tuned backbone. \\
\texttt{rag\_max\_context\_tokens}     & $384$ & design & Sets the input-length limit of the retrieval prompt ($8 \times 384 + 512$ tokens). \\
\texttt{rag\_max\_new\_tokens}         & $512$ & design & Generation budget under retrieval-augmented prompt. \\
\texttt{rag\_do\_sample}              & \texttt{false} & design & Greedy decoding for reproducibility. \\
\texttt{rag\_index\_type}             & \texttt{"ivf\_pq"} & design & FAISS index backend. \\
\texttt{rag\_faiss\_nlist},\allowbreak{} \texttt{nprobe\_tier3},\allowbreak{} \texttt{nprobe\_tier2} & $32768$, $64$, $16$ & design & IVF-PQ density and adaptive probe budget. \\
\texttt{rag\_top\_k}                  & $5$ & literature & Karpukhin et al.\ 2020 DPR standard. \\
\midrule
\multicolumn{4}{@{}l}{\textit{Memory store (\Cref{sec:memory-store})}} \\
\texttt{sbert\_model}                & \texttt{sentence-\allowbreak transformers/\allowbreak all-mpnet-\allowbreak base-v2} & literature & 768-dimensional sentence encoder. \\
\texttt{embedding\_dim}              & $768$ & literature & Must match the encoder output dimension. \\
\texttt{max\_memory\_size}            & $1\,000\,000$ & design & Capacity for thesis-scale and production-scale runs. \\
\texttt{memory\_index\_type}          & \texttt{"ivf\_pq"} & design & FAISS backend from $20{,}000$ entries; below that, and so throughout the run, search is exact. \\
\texttt{novelty\_threshold}          & $0.95$ & design & Skip storage if nearest-neighbour cosine similarity exceeds this. \\
\texttt{pruning\_trigger}            & $0.95$ & design & Trigger pruning when memory reaches this fraction of capacity (never reached). \\
\texttt{pruning\_amount}             & $0.20$ & design & Remove the bottom fraction by value score. \\
\texttt{value\_importance\_weight},\allowbreak{} \texttt{value\_recency\_weight} & $0.70$, $0.30$ & design & Value-score decomposition. \\
\midrule
\multicolumn{4}{@{}l}{\textit{Experiment scope}} \\
\texttt{num\_cycles}                 & $10$  & pre-specified & Registered horizon; the run stopped after cycle five, as planned after the cycle-four rollback. \\
\texttt{calibration\_set\_size}       & $500$ & design & Cycle-zero calibration-fold size per training benchmark. \\
\texttt{purity\_validation\_set\_size} & $500$ & design & Cycle-zero purity-validation slice per training benchmark. \\
\texttt{questions\_per\_cycle}        & $5\,000$ & design & Stream-chunk question budget per cycle. \\
\end{longtable}
}

Only four groups of quantities are refitted at each kept cycle boundary: the pooled and per-benchmark pre-routing temperatures, the per-benchmark safety floors, the per-signal isotonic curves, and the boost weights and intercept (\Cref{tab:calibration-trajectory}); every value in the registry is fixed for the whole run. The model and hardware are given in \Cref{app:repro-hardware,app:repro-versions}, and the cycle-zero validation gate in \Cref{sec:adj-gate}.

The retrieval-utility classifier (\Cref{sec:ruc}) has no fields in \verb|caem/config.py|. Its operating point lives in \verb|caem/ruc/v2/feature_spec.json|: the threshold $\tau_{\text{RUC}} = 0.375$, the regularisation strength $C = 0.5$, the thirty-six-feature ordering, the standard-scaler and logistic-head joblib paths, and the model-type guard \codepath{"logistic\_regression\_scaled"}. The classifier is loaded once at pipeline construction from the on-disk specification and is not refit at any cycle boundary; the architectural rationale is documented in \Cref{sec:ruc}.

\section{Runbook step graph}
\label{app:runbook}

The end-to-end execution sequence partitions into a pre-run setup block, a six-step main cycle loop, and a post-run analysis block. The pre-run block runs once before the first cycle; the main cycle loop repeats per cycle; the post-run block runs once after the final cycle closes. Paths: \verb|archive/scripts_superseded/old_runners/run_phase1a.sh| (the top-level driver used for the run, now archived), \verb|scripts/run_experiment.py| (the in-cycle implementation), and \verb|docs/PRODUCTION_RUNBOOK.md| Section~5 (the production version).

\paragraph{Pre-run setup (Steps 1--6):}
\begin{enumerate}
\item \textbf{Prompt smoke test} (\verb|step_5_9_prompt_smoke|). Renders $5$ samples per benchmark against the registered prompt templates and asserts the answer-line contract. Gate: every sample renders with a non-empty answer line.
\item \textbf{Alias dictionary build} (\verb|step_5_9_5_alias_dict|). Builds \verb|data/alias_dict.json| from the benchmarks' gold aliases. The alias signal is not part of the deployed composite; the dictionary is kept for ablation only.
\item \textbf{Cold-start seeding} (\verb|step_6_reseed| and \verb|step_6_5_em_prune|). Answers questions from each training benchmark's seed pool under a cold-start storage threshold of $0.50$ until about $200$ per benchmark are stored ($603$ answers), then removes the stored answers that are wrong against the gold labels, leaving $431$ entries (\Cref{sec:memory-store}).
\item \textbf{Cycle-zero evaluation and calibration-fold scoring} (\verb|step_7_0_cycle0|). Runs the pristine generator against the held-out evaluation fold and the calibration fold ($500$ samples per training benchmark, $1{,}500$ pooled). Writes per-sample nine-signal vectors and ground-truth labels. Gate: every JSON has the expected sample count and no NaN signals.
\item \textbf{Composite calibration and validation gate} (\codepath{step\_7\_0\_calibrate}, \codepath{step\_7\_0\_2\_5\_rescore\_eval}, \codepath{step\_7\_0\_3\_validate\_weights}). Fits the per-signal isotonic curves, the per-benchmark shrinkage-blended children, the boost-layer logistic regression, the per-benchmark temperature scalars, and the per-benchmark safety floors. Runs the cycle-zero validation gate (\Cref{sec:adj-gate}). Gate: cycle-zero canonical composite locked at \verb|outputs/cycle_0/composite_calibration.json|; never overwritten by later cycles.
\item \textbf{Retention baseline freeze and snapshot upload} (\verb|step_19_2_slice| and \verb|step_gdrive_upload_pre_main|). Records the pristine-model probe accuracies on MMLU, TriviaQA-test, and CommonsenseQA-test as the cycle-zero anchor for the $0.93$ retention floor. Uploads the pre-main snapshot so that a lost machine can be recovered from it (\Cref{app:repro-snapshot}).
\end{enumerate}

\paragraph{Main cycle loop (Steps 1--6, repeating per cycle):}
\begin{enumerate}
\item \textbf{Self-improvement fine-tune:} Draws the SIL training pool from memory entries above $\tau_{\text{train}} = 0.70$, applies the pool-reweighting layer, runs the LoRA fine-tune for three epochs, and probes the retention panel. Gate: every probe ratio against the cycle-zero baseline clears $0.93$. On failure, the adapter is rolled back to the previous cycle's checkpoint, the memory is preserved, and the recalibration and re-verification of step 2 are skipped under the asymmetric-rollback contract.
\item \textbf{Cycle-boundary recalibration:} Re-scores the calibration fold under the post-fine-tune model, refits the per-benchmark temperature scalars, refits the per-signal isotonic curves under the shrinkage prior, and runs the retroactive re-verification pass against the stored memory under the recalibrated verifier. Skipped on aborted cycles.
\item \textbf{Stream-chunk pre-routing pass:} Runs the per-query pipeline on the cycle's stream chunk ($2{,}000$ FEVER, $2{,}000$ TriviaQA, $700$ CommonsenseQA samples), accumulating storage admissions, deferred entries, abstentions, and discards.
\item \textbf{Stream-chunk verification and storage:} Applies the calibrated probability composite, the four-outcome decision tree, and the storage admission. Writes \verb|memory_store_cycle_N.{faiss,meta}| and \verb|deferred_buffer_cycle_N.pkl|.
\item \textbf{Held-out evaluation:} Runs the per-query pipeline on the held-out evaluation fold ($300$ samples per benchmark across the five-benchmark panel) under \verb|store_to_memory=False| so the eval fold does not contaminate the memory pool.
\item \textbf{Checkpoint and halt-trigger evaluation:} Writes the cycle's coverage diagnostic, the per-tier hit-rate roll-up, and the per-cycle composite calibration snapshot. Evaluates the automatic early-stop check.
\end{enumerate}

\paragraph{Post-run analysis:} the baseline re-run under the comparison protocol and the rebuild of every reported number; the two commands are given in \Cref{app:repro-steps}.

\paragraph{Asymmetric-rollback contract:} triggered when any retention probe drops below $0.93$ of its cycle-zero value in step 1 of the main cycle loop. The adapter reverts to the previous kept cycle's checkpoint, but the memory does not: answers stored during the failed cycle are kept. Step 2 (recalibration, retroactive re-verification and deferred reconsideration) is skipped, because the weights it would re-score are unchanged, and the re-verification log records the skip. Steps 3 to 6 still run, so memory grows and the rolled-back model is evaluated. In the run the contract fired once, at cycle four, when the TriviaQA-test probe ratio fell to $0.886$; four of the five fine-tunes were kept.
